\section{安装与编译}

\subsection{文件组成}

本宏包由两个文件组成：
\begin{itemize}
  \item \tkzname{tkz-fn.sty}：\LaTeX{} 接口（宏和选项）；
  \item \tkzname{tkz\_fn.lua}：Lua 代码（计算和生成的 \TIKZ\ 代码）。
\end{itemize}

这两个文件必须放置在 \TeX{} 可识别的文件夹中，例如 \tkzname{texmf/tex/latex/tkz-fn}，或放在文档所在文件夹中。

\subsection{加载方式}

\begin{tkzltxexample}[]
\usepackage{tkz-fn}         % 格式随 babel 语言自动选择
\usepackage[us]{tkz-fn}     % 结果使用美式格式
\usepackage[fr]{tkz-fn}     % 结果使用法式格式
\end{tkzltxexample}

\TIKZ\ 及其 \tkzname{plotmarks} 库会自动加载。示例中用到的某些 \TIKZ\ 选项需要其他库：|\usetikzlibrary{arrows,patterns}|。

\subsection{编译方式}

文档必须使用 \tkzname{LuaLaTeX} 编译。不需要 \tkzname{shell-escape} 选项，也不会生成辅助文件。

\subsection{函数的书写}

表达式即为 Lua 表达式。函数的变量为 |x|，参数曲线或极坐标曲线的变量为 |t|。

\begin{center}
\begin{tabular}{lll}
\toprule
Lua & 旧语法 (gnuplot) & 说明 \\
\midrule
|x^3|, |x^0.5| & |x**3|, |x**0.5| & 幂运算 \\
|1/3| & |1./3| & 除法（始终为小数） \\
|ln(x)| 或 |log(x)| & |log(x)| & 自然对数 \\
|log10(x)| & |log10(x)| & 常用对数 \\
|exp(x)|, |e| & |exp(x)|, |exp(1)| & 指数函数，数 $\text{e}$ \\
|sqrt(x)|, |abs(x)| & |sqrt(x)|, |abs(x)| & 平方根，绝对值 \\
|sin|, |cos|, |tan| & 相同 & 弧度制 \\
|asin|, |acos|, |atan| & 相同 & 反三角函数 \\
|sinh|, |cosh|, |tanh| & 相同 & 双曲函数 \\
|pi| & |pi| & $\pi$ \\
|floor|, |ceil|, |sign|, |mod(a,b)| & |floor|, |ceil|, |sgn|, |%| & \\
|cbrt(x)|, |root(x,n)| & & 立方根，$n$ 次方根 \\
\bottomrule
\end{tabular}
\end{center}

\tkzHand 用 \tkzcname{tkzFnDef} 定义的函数可在其他表达式中使用：|2*f(x)+1|。全局 Lua 变量可用作参数：|\directlua{a = 0.5}| 然后写 |a*x^2|。

\tkzHand 函数 |diff(f,x)|、|integral(f,a,b)| 和 |riemann(f,a,b,n,"left")| 分别给出某点处的导数、积分（辛普森法则）和黎曼和。它们与 \tkzcname{tkzFnPrint} 配合使用。

\tkzBomb 圆括号内的参数不能包含未保护的右括号：应写 |(f,{1/exp(1)})|，而非 |(f,1/exp(1))|。

\endinput
